\bigskip
% lecture 4, page 1

\begin{center}
\fbox{\textbf{ALGEBRA I}}\\[2pt]
(Lecture 4)
\end{center}
\begin{flushright}
\textit{20 Oct.\ 1994}
\end{flushright}

\bigskip
\noindent\textbf{(1)} Composition law induced on a stable subset. $(A,\varphi)$, $\varphi : A\times A \to A$, $A\neq\varnothing$.

\noindent $B\subseteq A$, $B$ a stable subset of $A$ with respect to $\varphi$ if $(\forall)\,x,y\in B \Rightarrow \varphi(x,y)\in B$.

\noindent $\varphi' : B\times B \to B$.

\bigskip
\noindent\textbf{(2)} Composition law induced on a factor set. $(A,R)$

\noindent $\varphi : A\times A \to A$, \quad $R\subseteq A\times A$

\noindent $R$ compatible with $\varphi$: \quad $xRx'$ and $yRy'$ $\Rightarrow$ $xyRx'y'$ (in multiplicative notation.)

\[
\frac{A}{R} = \left\{\hat a \mid a\in A\right\} \qquad \varphi^* : \frac{A}{R}\times\frac{A}{R} \to \frac{A}{R} \qquad \varphi^*(\hat a,\hat b) \overset{\text{def}}{=} \widehat{\varphi(a,b)}
\]
\[
\hat a\hat b = \widehat{ab}; \qquad \hat a + \hat b = \widehat{a+b}.
\]

\bigskip
\noindent\textbf{(3)} Composition law induced on a set of functions. $(A,\varphi)$, $I\neq\varnothing$.

\noindent $A^I = \{f \mid f:I\to A\}$

\noindent $\tilde\varphi$ = the composition law induced by $\varphi$.

\noindent $(A^I,\tilde\varphi)$ \qquad $(\mathbb{R},+)$; $I=[a,b]$.

\noindent $\mathbb{R}^{[a,b]} = \{f \mid f:[a,b]\to\mathbb{R}\}$; \qquad $\mathbb{R}^{\{1,\ldots,n\}} = \{f \mid f:\{1,\ldots,n\}\to\mathbb{R}\}$.

\medskip
\noindent $A = a_0,a_1,\ldots,a_n,\ldots$

\noindent $m,n\in\mathbb{N}^*$.

\[
\mathbb{R}^{\{1,\ldots,m\}\times\{1,\ldots,n\}}
\]
\noindent $A : \{1,\ldots,m\}\times\{1,\ldots,n\} \to \mathbb{R}$ \hfill $\mathbb{R}\ni a_{ij} = A(i,j)$

\medskip
\noindent $(A,\varphi)$, $I\neq\varnothing$

\noindent $f,g\in A^I$

\[
f*g : I\to A \qquad (f*g)(x) \overset{\text{def}}{=} f(x)*g(x) \in A.
\]
\[
\varphi^* : A^I\times A^I \to A^I \qquad \big(\varphi^*(f,g)\big)(x) \overset{\text{def}}{=} \varphi\big(f(x),g(x)\big).
\]

\medskip
\noindent (1) Let $f,g : [a,b]\to\mathbb{R}$; \quad $f+g : [a,b]\to\mathbb{R}$.

\noindent $(\forall)\, x\in[a,b]$ \qquad $(f+g)(x) = f(x)+g(x)$.

\noindent $fg : [a,b]\to\mathbb{R}$ \qquad $(fg)(x) = f(x)g(x)$.




% lecture 4, page 2
\section*{Lecture 4 - Page 2}

\noindent\fbox{Theorem} \ Let $(A,\varphi)$; $\varphi : A\times A \to A$, $I\neq\varnothing$.

\[
\tilde\varphi : A^I \times A^I \to A^I
\]

\begin{enumerate}
\item[(1)] If $\varphi$ is associative (commutative) $\Rightarrow$ $\tilde\varphi$ is associative (commutative).
\item[(2)] If $e\in A$ is the identity element for $\varphi$ $\Rightarrow$ $\tilde e : I\to A$, $\tilde e(i)=e$, is the identity element for $\tilde\varphi$.
\item[(3)] $\varphi$ has an identity element and $f(i)$ is invertible with respect to $\varphi$ for every $i$ $\Rightarrow$ $f$ is invertible with respect to $\tilde\varphi$, and an inverse of $f$ is $f'(i) = (f(i))'$.
\end{enumerate}

\bigskip
\noindent (1) We denote the law on $A$ by $*$.\footnote{the exact wording is unclear in the original (compressed/hard to read) -- it appears the law $\varphi$ is relabelled $*$ for convenience in the proof.}

\[
(f*g)(i) \overset{\text{def}}{=} f(i)*g(i) = g(i)*f(i) \overset{\text{def}}{=} (g*f)(i), \quad \text{true.}
\]

\noindent (2) $(f*\tilde e)(i) = f(i)*\tilde e(i) = f(i)$.

\noindent $f*\tilde e = \tilde e*f = f$.

\noindent (3) $(f*f')(i) = \tilde e(i) = e$.

\noindent $f(i)*f'(i) = f(i)*(f(i))' = e$.

\bigskip
\noindent\textbf{(4)} Product of two composition laws.

\noindent $(A_1,\varphi_1)$, $(A_2,\varphi_2)$; $A = A_1\times A_2$.

\noindent $\varphi : A\times A \to A$; \quad $\varphi : (A_1\times A_2)\times(A_1\times A_2) \to A_1\times A_2$

\[
\varphi\big((x_1,x_2),(y_1,y_2)\big) \overset{\text{def}}{=} \big(\varphi_1(x_1,y_1),\varphi_2(x_2,y_2)\big).
\]
\[
(x_1,x_2)+(y_1,y_2) = (x_1+y_1,\,x_2+y_2) \quad \text{-- additive notation.}
\]

\noindent $(\mathbb{Z},+)$ \qquad $(\mathbb{Z}_6,+)$ \qquad $\mathbb{Z}\times\mathbb{Z}_6$

\bigskip
\noindent\fbox{Theorem} \ Let $(A_1,\varphi_1)$, $(A_2,\varphi_2)$ and $\varphi = \varphi_1\times\varphi_2$.

\begin{enumerate}
\item[(1)] If $\varphi_1$ and $\varphi_2$ are associative (commutative) $\Rightarrow$ $\varphi$ is associative (commutative)
\item[(2)] $e_1,e_2$ the identity elements for $\varphi_1,\varphi_2$ $\Rightarrow$ $(e_1,e_2)$ is the identity element for $\varphi$.
\item[(3)] $x_1\in A_1$ invertible with respect to $\varphi_1$ and $x_2\in A_2$ invertible with respect to $\varphi_2$ $\Rightarrow$ $(x_1,x_2)$ is invertible with respect to $\varphi$, and $(x_1,x_2)' = (x_1',x_2')$.
\end{enumerate}




% lecture 4, page 3
\section*{Lecture 4 - Page 3}

\noindent\textbf{(5)} Transport of a composition law through a bijection

\[
(A,\varphi) \qquad A \xrightarrow{f} A'
\]

\noindent If $x',y'\in A' \Rightarrow \exists!\, x,y\in A$ such that $x'=f(x)$ and $y'=f(y)$

\[
\varphi' : A'\times A' \to A', \qquad \varphi'(x',y') = f\big(\varphi(x,y)\big) \in A'
\]

\noindent [notation: $\varphi(x,y)=xy$] \quad $x'y' \overset{\text{def}}{=} f(xy)$\footnote{this line is very compressed in the original; it appears to introduce the multiplicative convention $\varphi(x,y)=xy$, so that the definition of $\varphi'$ becomes $x'y'=f(xy)$.}

\[
f : A\to A' \qquad f(xy) = f(x)f(y) \quad (\forall)\,x,y\in A.
\]
\[
f^{-1}(x'y') = f^{-1}(x')f^{-1}(y') \quad (\forall)\,x',y'\in A'.
\]

\bigskip
\noindent\underline{Theorem:} Let $(A,\varphi)$, $\varphi:A\times A\to A$.

\noindent $f:A\to A'$ and $\varphi':A'\times A'\to A'$.

\begin{enumerate}
\item[1)] $\varphi$ associative (commutative) $\Leftrightarrow$ $\varphi'$ is associative (commutative).
\item[2)] $e\in A$ the identity element for $\varphi$ $\Leftrightarrow$ $f(e)$ is the identity element for $\varphi'$.\footnote{the original reads ``for $f$'' / ``for $f'$'' -- corrected here to $\varphi$/$\varphi'$ for consistency with point 1) and the rest of the theorem.}
\item[3)] $x\in A$ invertible with respect to $\varphi$ $\Leftrightarrow$ $f(x)$ invertible with respect to $\varphi'$.
\end{enumerate}

\noindent $x',y'\in A'$.

\[
x'y' = f(x)f(y) = f(xy) = f(yx) = f(y)f(x).
\]

\bigskip
\begin{center}
\fbox{\Large\textbf{GROUPS}}
\end{center}

\noindent\underline{\textbf{Subgroups. Lagrange's Theorem}}

\bigskip
\noindent\underline{Def:} Let $G$ be a group. A nonempty subset $H$ of $G$ is called a \textbf{subgroup} of the group $G$ if:

\begin{enumerate}
\item[(1)] $(\forall)\,x,y\in H \Rightarrow xy\in H$
\item[(2)] $(\forall)\,x\in H \Rightarrow x^{-1}\in H$
\end{enumerate}

\bigskip
\noindent\underline{Problem:} Given a group $G$, determine all congruences on $G$.

\noindent An equivalence relation on $G$ compatible with the group operation:

\[
xRy \Rightarrow zxRzy.
\]



% lecture 4, page 4
\section*{Lecture 4 - Page 4}

\noindent $n\in\mathbb{N}$.

\bigskip
\noindent $H \overset{\text{not}}{=} n\mathbb{Z} = \{ny \mid y\in\mathbb{Z}\}$ is a subgroup of the group $(\mathbb{Z},+)$.

\noindent $\Rightarrow \exists!\, n\in\mathbb{N}$ such that $H=n\mathbb{Z}$.

\begin{enumerate}
\item[1)] $(\forall)\, x,y\in n\mathbb{Z} \Rightarrow x+y\in n\mathbb{Z}$.
\item[2)] $(\forall)\, x\in n\mathbb{Z} \Rightarrow -x\in n\mathbb{Z}$.
\end{enumerate}

\[
\left.\begin{aligned}
x&=nq_1\\
y&=nq_2
\end{aligned}\right\} \Rightarrow x+y = n(q_1+q_2).
\]

\bigskip
\noindent\underline{Converse.}

\noindent $H$ is a subgroup of $(\mathbb{Z},+)$ $\Rightarrow$ $H=n\mathbb{Z}$.

\noindent $0\in H$

\noindent $H=\{0\}$: \quad $\{0\} = \{0y \mid y\in\mathbb{Z}\}$.

\noindent Suppose $H\neq\{0\}$.

\[
H^+ = \{x\in H \mid x>0\}.
\]

\noindent Let $n$ be the smallest element of $H^+$

\noindent We show that $H=n\mathbb{Z}$.

\bigskip
\noindent (by induction)\footnote{the exact word/symbol is unclear in the original -- it appears to indicate an inductive argument.} $x=nq$.

\noindent $q=1 \Rightarrow x=n$

\noindent $nk\in H \Rightarrow n(k+1)\in H \Rightarrow nk+n\in H$.

\noindent Hence $n\mathbb{Z}\subseteq H$

\bigskip
\noindent Suppose\footnote{``Pp.'' could also mean ``by the proposition [of division with remainder]''.} $x=nq+r$.

\noindent Suppose $r\neq 0 \Rightarrow r=x-nq\in H$.

\bigskip
\noindent\underline{\textbf{Left (right) congruences on a group.}}

\noindent\underline{\textbf{Lagrange's Theorem.}}

\bigskip
\noindent Let $G$ be a group, $H\subseteq G$

\[
R_H^L = \{(x,y)\in G\times G \mid x^{-1}y\in H\}.
\]
\[
R_H^R = \{(x,y)\in G\times G \mid xy^{-1}\in H\}.
\]



% lecture 4, page 5
\section*{Lecture 4 - Page 5}

\noindent $xR_H^L y \Leftrightarrow (x,y)\in R_H^L$.

\noindent $x\equiv_L y \pmod H \Leftrightarrow x^{-1}y\in H$.

\bigskip
\noindent $x\equiv_L y \pmod H \Leftrightarrow y\equiv_L x \pmod H$ $(\forall)\, x,y\in G$.

\noindent $x\equiv_L y \pmod H$, $y\equiv_L z \pmod H \Rightarrow x\equiv_L z \pmod H$

\bigskip
\noindent $x\equiv_L y \pmod H \Rightarrow zx\equiv_L zy \pmod H$. \ $(\forall)\, x,y\in G$.\footnote{``z'' is unclear in the original (it could look like the digit ``2''), but the meaning matches the left-compatibility established on the previous page ($xRy\Rightarrow zxRzy$).}

\bigskip
\noindent\underline{Remark}: If $R$ is an equivalence relation on $G$ compatible with left multiplication $\Rightarrow$ $\exists\, H\leq G$ such that $R=R_H^L$

\[
H \overset{\text{def}}{=} [e]_R = \{x\in G \mid xRe\}.
\]

\noindent $x,y\in H \overset{?}{\Rightarrow} xy\in H$.

\noindent $xRe$

\noindent $yRe$

\noindent $\Rightarrow xyRe \Rightarrow xy\in H$.

\bigskip
\noindent We show that $R=R_H^L$.

\noindent Let $(x,y)\in R \Leftrightarrow xRy \Leftrightarrow eRx^{-1}y \Leftrightarrow x^{-1}yRe \Leftrightarrow x^{-1}y\in H$

\noindent $\Leftrightarrow (x,y)\in R_H^L$.

\bigskip
\noindent\underline{Prop}: $G$ a group and $H\leq G$; and $a\in G$ such that

\[
[a]_{R_H^L} = aH \overset{\text{def}}{=} \{ah \mid h\in H\}
\]

\noindent = the left coset with representative $a$ modulo the subgroup $H$.

\bigskip
\noindent Suppose $x\in[a]_{R_H^L} \Rightarrow x\equiv_L a \pmod H \Rightarrow a\equiv_L x \pmod H$

\noindent $\Leftrightarrow a^{-1}x\in H \Leftrightarrow x=ah,\ h\in H$.

\bigskip
\noindent Let $\{a_i\}_{i\in I}$ be a right transversal of $H$ in $G$.

\[
\{a_i\}_{i\in I} \xrightarrow{\ \xi\ } \{a_i^{-1}\}_{i\in I} \qquad
\begin{aligned}
&(\forall)\, i\neq j \Rightarrow a_i \not\equiv_L a_j \pmod H\\
&(\forall)\, x\in G,\ \exists!\, i\in I \text{ such that } x\equiv_L a_i \pmod H.
\end{aligned}
\]

\noindent $\xi$ is a bijection.

\[
\text{Card}(I) = [G:H]
\]

\noindent = the index of $H$ in $G$. (the number of left (right) residue classes modulo the subgroup $H$).



% lecture 4, page 6
\section*{Lecture 4 - Page 6}

\noindent\underline{Proof}:

\bigskip
\noindent (1) Suppose $a_i^{-1}\equiv_R a_j^{-1} \pmod H \Rightarrow a_i^{-1}a_j\in H \Rightarrow a_i\equiv_L a_j \pmod H$ \ $\bot$.

\bigskip
\noindent (2) $x\in G \Rightarrow \exists\, i\in I$ such that $x^{-1}\equiv_L a_i \pmod H \Rightarrow (x^{-1})^{-1}a_i\in H \Rightarrow$

\noindent and $xa_i\in H \Rightarrow x(a_i^{-1})^{-1}\in H \Rightarrow x\equiv_R a_i^{-1} \pmod H$.

\bigskip
\noindent $G$ a group, \ $f:G\to G$, \ $f(x)=x^{-1}$ is a bijection.

\bigskip
\noindent
\[
G = \bigcup_{i\in I} a_iH
\]

\noindent $|G|<\infty \Rightarrow |I|=d<\infty$

\medskip
\noindent Suppose $a_1,a_2,\ldots,a_d$ is a left transversal of $H$ in $G$.

\[
|G|=n \qquad |H|=m
\]

\noindent $|G| = |H|[G:H]$ -- Lagrange's Theorem

\noindent $|H|=|aH|$.

\[
|G| = \sum_{i=1}^{d} |a_iH| = \underbrace{|H|+\cdots+|H|}_{d} = d|H| = dm.
\]

\noindent $f : H \to aH$

\noindent $f$ is surjective.

\noindent $f(x)=ax$; \quad $ax_1=ax_2 \Rightarrow x_1=x_2$ \quad $\Rightarrow$ \ $f$ is bijective.

\bigskip
\noindent\underline{Ex}: a group with 12 elements that admits a subgroup of 4 elements.

\[
(\mathbb{Z}_{12},+)
\]
\[
\hat 0,\hat 1,\hat 2,\hat 3,\hat 4,\hat 5,\hat 6,\hat 7,\hat 8,\hat 9,\hat{10},\hat{11}
\]
\[
\{\hat 0,\hat 3,\hat 6,\hat 9\} \leq (\mathbb{Z}_{12},+).
\]

